\chapter*{Conclusion}
\addcontentsline{toc}{chapter}{Conclusion}

\markboth{Conclusion}{Conclusion}

Ce projet a permis de concevoir et de développer le frontend d'une application mobile de réseau social complète, \textbf{Sharely}, en adoptant les technologies et les pratiques les plus en vogue dans l'industrie du développement mobile.

\section*{Bilan Technique}

L'architecture retenue --- React Native avec Expo SDK~54, Expo Router pour la navigation, et l'API Context pour la gestion d'état --- a prouvé son efficacité pour un projet de cette envergure. Le patron deux fichiers (\texttt{app/} $\to$ \texttt{components/}) a maintenu une séparation nette entre la déclaration des routes et la logique métier. Le recours aux \texttt{refs} React pour les données non-réactives (curseurs, pools), l'Optimistic UI pour les interactions, et les connexions WebSocket en singleton ont permis d'atteindre un niveau de réactivité et de fluidité satisfaisant.

Le système de thèmes via des tokens de couleurs sémantiques a garanti une cohérence visuelle entre le mode clair et le mode sombre à travers toute l'interface, tandis que les animations ciblées (J'aime, progress bar de story, modals glissants) ont enrichi l'expérience utilisateur sans compromettre les performances.

\section*{Fonctionnalités Réalisées}

L'application couvre un spectre fonctionnel large :
\begin{itemize}[leftmargin=1.5em, itemsep=0.2em]
    \item Authentification complète avec persistance de session et renouvellement automatique.
    \item Fil d'actualité infini avec injection de publicités et mode découverte.
    \item Stories avec filtres, superpositions interactives et musique.
    \item Reels en défilement vertical avec lecture automatique.
    \item Messagerie directe en temps réel via WebSockets.
    \item Exploration avec recherche d'utilisateurs et grille de découverte.
    \item Notifications en temps réel sous forme de toasts et panneau dédié.
    \item Chatbot IA intégré accessible depuis tous les écrans.
\end{itemize}

\section*{Perspectives d'Amélioration}

Bien que l'application soit fonctionnelle, plusieurs axes d'amélioration méritent d'être envisagés dans des développements futurs :

\begin{description}[leftmargin=1.8em]
    \item[\textbf{Hooks personnalisés.}] Extraire la logique d'état et les appels API des composants vers des hooks dédiés (\texttt{useInfiniteFeed}, \texttt{useSocket}, \texttt{useProfile}) améliorerait la testabilité et la réutilisabilité.

    \item[\textbf{Tests automatisés.}] L'absence de tests unitaires et d'intégration représente un risque à mesure que la base de code grandit. L'adoption de Jest avec Testing Library permettrait de couvrir les composants et la logique métier.

    \item[\textbf{Optimisation des images.}] L'intégration d'une bibliothèque comme \texttt{expo-image} avec mise en cache agressive réduirait la consommation de données et améliorerait les temps de chargement.

    \item[\textbf{Internationalisation.}] L'ajout d'un support multilingue (Arabic, Français, Anglais) via \texttt{i18n-js} ou \texttt{react-intl} élargirait l'audience cible de l'application.

    \item[\textbf{Accessibilité.}] Un audit d'accessibilité (propriétés \texttt{accessibilityLabel}, \texttt{accessibilityRole}) rendrait l'application utilisable par les personnes en situation de handicap.
\end{description}

\section*{Conclusion Générale}

Ce projet a été une occasion unique de mettre en pratique, dans un contexte réaliste, l'ensemble des compétences abordées dans le module Développement Mobile. La complexité fonctionnelle de Sharely --- combinant navigation avancée, temps réel, multimédia et thèmes dynamiques --- a révélé la puissance de l'écosystème React Native / Expo pour développer des applications mobiles de qualité professionnelle. Il constitue une base solide sur laquelle pourront s'appuyer des évolutions futures aussi bien techniques que fonctionnelles.
